% --- French Abstract ---

%\chapter*{Résumé}
\noindent{\Large\bfseries Résumé}\par\vspace{12pt}
\thispagestyle{empty}
\enlargethispage{\baselineskip}
%\vspace*{0.5cm}
\noindent Ce rapport présente le travail réalisé dans le cadre de notre projet de fin d'études pour l'obtention de la Licence Nationale en Informatique de Gestion, spécialité Business Intelligence, à l'Institut Supérieur de Gestion Industrielle de Sfax (ISGIS). Effectué au sein de la Fondation Djagora, ce projet nous a permis de contribuer à la conception et au développement de LogiLink, une plateforme intelligente basée sur l'intelligence artificielle et dédiée à l'adéquation formation-emploi pour les étudiants de l'ISGIS. Cette solution vise à répondre aux enjeux de l'employabilité dans les secteurs du transport et de la logistique en facilitant l'orientation professionnelle et l'insertion sur le marché du travail. Dans ce cadre, notre mission s'est principalement concentrée sur le développement de fonctionnalités stratégiques, notamment la génération automatisée de CV, la mise en place de mécanismes d'évaluation fondés sur le calcul du score ATS et du score d'employabilité, ainsi que le développement d'une fonctionnalité permettant de prévoir l'impact des certifications sur l'employabilité des étudiants à travers le calcul d'un score d'employabilité prévisionnel.

\vspace{6pt}

\noindent \textbf{Mots-clés :} Transport, Logistique, Génération automatisée de CV, Score ATS, Score d'employabilité, Certifications, Score d'employabilité prévisionnel.

\vspace{0.5cm}

\noindent{\Large\bfseries Abstract}\par\vspace{12pt}

\noindent This report presents the work carried out as part of our final year project in order to obtain the National Bachelor's Degree in Business Computing, specializing in Business Intelligence, at the Higher Institute of Industrial Management of Sfax (ISGIS). Conducted within the Djagora Foundation, this project allowed us to contribute to the design and development of LogiLink, an artificial intelligence-based platform dedicated to aligning education and employment for ISGIS students. This solution aims to address employability challenges in the transport and logistics sectors by facilitating career guidance and professional integration into the job market. In this context, our work mainly focused on the development of key functionalities, including automated CV generation, the implementation of evaluation mechanisms based on ATS scoring and employability scoring, as well as the development of a feature designed to predict the impact of certifications on students' employability through a predictive employability score.

\vspace{6pt}

\noindent \textbf{Keywords:} Transport, Logistics, Automated CV generation, ATS scoring, Employability scoring, Certifications, Predictive employability score.